\documentclass[12pt]{article}
\usepackage[utf8]{inputenc}
\usepackage[T1]{fontenc}
\usepackage{hyperref}
\title{Tetete Language Resources for K-12 Curriculum Development}
\author{Compiled by Claude Sonnet 5 (claude-sonnet-5) for Andr\'as Kornai}
\date{}
\begin{document}
\maketitle
\section{Introduction}
Tetete (Tetet\'e, ISO 639-3 teb, Glottocode tete1252) is an extinct Western Tucanoan language of the Ecuador--Colombia border area, reported as close to Secoya/Siona. The prompt's family label ``Bookkeeping'' appears to be a data error; sources found describe it as Tucanoan. The last known survivors were three elderly people encountered in 1966, and the last outside contact was in 1973. No native-language teaching materials were found in search.
\section{Primary Educational Resources}
None found. No primers, textbooks, or word lists for Tetete turned up in search.
\section{Linguistic and Reference Materials}
Encyclopedic overviews only: \url{https://en.wikipedia.org/wiki/Tetet%C3%A9_language} and \url{https://endangeredlanguages.com/lang/2151}. A related catalogue page: \url{https://rosettapanglossia.longnow.org/index.php?title=Teb}.
\section{Religious and Community Resources}
No community organisation or scripture was found. A page on the history of the Tetete people (including the contested account of their extermination): \url{https://en.wikipedia.org/wiki/Tetet%C3%A9_people}.
\section{Digital and Audio Resources}
None found.
\section{Curriculum Development Context}
With no documented lexicon or grammar located, a Tetete-medium curriculum is not feasible from what was found. A heritage curriculum could be history-focused, and could draw on the closely related, living Siona/Secoya languages for comparative material (not searched here).
\end{document}